% GENERATED FILE. Edit canonical lessons, then run npm run generate:books.
% canonical-source-hash: fnv1a64:25af0652d521cc4a
% canonical-lessons: IT-C279-interruzione, IT-C279-sospendere, IT-C279-rimandare, IT-C279-in-anticipo, IT-C279-tappa

\chapter{Più tardi o prima}
\label{ch:it-piu-tardi-o-prima}
\hlchaptermodality{279}

\begin{chapteropening}
\textbf{By the end of this chapter:} \emph{I can say that something was interrupted, suspended, put off or done early, and name a stop along the way.}

Complete the last lesson of chapter 279: I can say that something was interrupted, suspended, put off or done early, and name a stop along the way.
\end{chapteropening}

\section[l'interruzione]{l'interruzione — the interruption, the break}

\label{lesson:IT-C279-interruzione}



\begin{quote}
Before the new one: say the Italian for the conclusion, then the Italian for to resume.
\end{quote}

\subsection*{You'll want to know: l'interruzione}
\textbf{l'interruzione} — \textquotedblleft{}the interruption, the break\textquotedblright{}.

\begin{quote}
C'è stata un'interruzione di corrente.
\end{quote}

\begin{quote}
\emph{There was a power cut.}
\end{quote}

\emph{Senza interruzione} — without a break.

\begin{cousinweb}
\emph{interruzione} is a learned borrowing from Latin \emph{interruptiō} \textquotedblleft{}a breaking off\textquotedblright{} (EtymDB), from \emph{interrumpere}, built on \emph{rumpere} \textquotedblleft{}to break\textquotedblright{}.
\end{cousinweb}

\begin{grammarlens}[title={What you've built}]
One more word for telling, explaining and describing at B1.
\end{grammarlens}

\subsection*{Your turn}
\begin{itemize}
\raggedright
  \item \emph{Say it:} \emph{l'interruzione}
  \item \emph{Say it:} \emph{l'interruzione}, once more
  \item \emph{Say it:} a sentence of your own with \emph{l'interruzione}
  \item \emph{Recall:} say the Italian for the conclusion, then the Italian for to resume, then say \emph{l'interruzione} again
\end{itemize}

\begin{tcolorbox}[breakable,colback=teal!4,colframe=teal!35!black,title={Before you move on}]
What does l'interruzione mean? (\textquotedblleft{}The interruption, the break\textquotedblright{}.) Say it once more.
\end{tcolorbox}
\section[sospendere]{sospendere — to suspend, to halt for a while}

\label{lesson:IT-C279-sospendere}



\begin{quote}
Before the new one: say the Italian for to resume, then the Italian for the interruption.
\end{quote}

\subsection*{You'll want to know: sospendere}
\textbf{sospendere} — \textquotedblleft{}to suspend, to halt for a while\textquotedblright{}.

\begin{quote}
Hanno sospeso la partita per la pioggia.
\end{quote}

\begin{quote}
\emph{They suspended the match because of the rain.}
\end{quote}

Past participle \emph{sospeso}. It also means ‘to hang up’: \emph{sospendere una lampada}.

\begin{cousinweb}
\emph{sospendere} is inherited from Latin \emph{suspendere} \textquotedblleft{}to hang up, to suspend\textquotedblright{} (EtymDB), from \emph{sub} \textquotedblleft{}under\textquotedblright{} and \emph{pendere} \textquotedblleft{}to hang\textquotedblright{}.
\end{cousinweb}

\begin{grammarlens}[title={What you've built}]
One more word for telling, explaining and describing at B1.
\end{grammarlens}

\subsection*{Your turn}
\begin{itemize}
\raggedright
  \item \emph{Say it:} \emph{sospendere}
  \item \emph{Say it:} \emph{sospendere}, once more
  \item \emph{Say it:} a sentence of your own with \emph{sospendere}
  \item \emph{Recall:} say the Italian for to resume, then the Italian for the interruption, then say \emph{sospendere} again
\end{itemize}

\begin{tcolorbox}[breakable,colback=teal!4,colframe=teal!35!black,title={Before you move on}]
What does sospendere mean? (\textquotedblleft{}To suspend, to halt for a while\textquotedblright{}.) Say it once more.
\end{tcolorbox}
\section[rimandare]{rimandare — to put off, to postpone; to send back}

\label{lesson:IT-C279-rimandare}



\begin{quote}
Before the new one: say the Italian for the interruption, then the Italian for to suspend.
\end{quote}

\subsection*{You'll want to know: rimandare}
\textbf{rimandare} — \textquotedblleft{}to put off, to postpone; to send back\textquotedblright{}.

\begin{quote}
Abbiamo rimandato la cena a sabato.
\end{quote}

\begin{quote}
\emph{We put the dinner off until Saturday.}
\end{quote}

The everyday word for putting something off: \emph{rimandare a domani}.

\begin{cousinweb}
\emph{rimandare} is \emph{ri-} \textquotedblleft{}again, back\textquotedblright{} with \emph{mandare} \textquotedblleft{}to send\textquotedblright{}, which is inherited from Latin \emph{mandāre} \textquotedblleft{}to hand over, to entrust\textquotedblright{} (Lewis \& Short).
\end{cousinweb}

\begin{grammarlens}[title={What you've built}]
One more word for telling, explaining and describing at B1.
\end{grammarlens}

\subsection*{Your turn}
\begin{itemize}
\raggedright
  \item \emph{Say it:} \emph{rimandare}
  \item \emph{Say it:} \emph{rimandare}, once more
  \item \emph{Say it:} a sentence of your own with \emph{rimandare}
  \item \emph{Recall:} say the Italian for the interruption, then the Italian for to suspend, then say \emph{rimandare} again
\end{itemize}

\begin{tcolorbox}[breakable,colback=teal!4,colframe=teal!35!black,title={Before you move on}]
What does rimandare mean? (\textquotedblleft{}To put off, to postpone; to send back\textquotedblright{}.) Say it once more.
\end{tcolorbox}
\section[in anticipo]{in anticipo — early, in advance}

\label{lesson:IT-C279-in-anticipo}



\begin{quote}
Before the new one: say the Italian for to suspend, then the Italian for to put off.
\end{quote}

\subsection*{You'll want to know: in anticipo}
\textbf{in anticipo} — \textquotedblleft{}early, in advance\textquotedblright{}.

\begin{quote}
Siamo arrivati in anticipo di dieci minuti.
\end{quote}

\begin{quote}
\emph{We arrived ten minutes early.}
\end{quote}

The opposite of \emph{in ritardo}.

\begin{cousinweb}
\emph{anticipo} is built on \emph{anticipare} \textquotedblleft{}to do early\textquotedblright{}, a learned borrowing from Latin \emph{anticipāre} \textquotedblleft{}to take beforehand\textquotedblright{}, from \emph{ante} \textquotedblleft{}before\textquotedblright{} and \emph{capere} \textquotedblleft{}to take\textquotedblright{} (Lewis \& Short).
\end{cousinweb}

\begin{grammarlens}[title={What you've built}]
One more word for telling, explaining and describing at B1.
\end{grammarlens}

\subsection*{Your turn}
\begin{itemize}
\raggedright
  \item \emph{Say it:} \emph{in anticipo}
  \item \emph{Say it:} \emph{in anticipo}, once more
  \item \emph{Say it:} a sentence of your own with \emph{in anticipo}
  \item \emph{Recall:} say the Italian for to suspend, then the Italian for to put off, then say \emph{in anticipo} again
\end{itemize}

\begin{tcolorbox}[breakable,colback=teal!4,colframe=teal!35!black,title={Before you move on}]
What does in anticipo mean? (\textquotedblleft{}Early, in advance\textquotedblright{}.) Say it once more.
\end{tcolorbox}
\section[la tappa]{la tappa — the stage, the stop (on a journey)}

\label{lesson:IT-C279-tappa}



\begin{quote}
Before the new one: say the Italian for to put off, then the Italian for early.
\end{quote}

\subsection*{You'll want to know: la tappa}
\textbf{la tappa} — \textquotedblleft{}the stage, the stop (on a journey)\textquotedblright{}.

\begin{quote}
La prima tappa è stata un piccolo paese di montagna.
\end{quote}

\begin{quote}
\emph{The first stop was a small mountain village.}
\end{quote}

A race or a tour goes \emph{tappa per tappa}, stage by stage.

\begin{cousinweb}
\emph{tappa} was borrowed from French \emph{étape} \textquotedblleft{}a stopping place\textquotedblright{} (EtymDB), which goes back to a Germanic word for a pile or a storehouse.
\end{cousinweb}

\begin{grammarlens}[title={What you've built}]
One more word for telling, explaining and describing at B1.
\end{grammarlens}

\subsection*{Your turn}
\begin{itemize}
\raggedright
  \item \emph{Say it:} \emph{la tappa}
  \item \emph{Say it:} \emph{la tappa}, once more
  \item \emph{Say it:} a sentence of your own with \emph{la tappa}
  \item \emph{Recall:} say the Italian for to put off, then the Italian for early, then say \emph{la tappa} again
\end{itemize}

\begin{tcolorbox}[breakable,colback=teal!4,colframe=teal!35!black,title={Before you move on}]
What does la tappa mean? (\textquotedblleft{}The stage, the stop (on a journey)\textquotedblright{}.) Say it once more.
\end{tcolorbox}
